\begin{abstract}
At watersheds with a legacy source, per- and polyfluoroalkyl substances (PFAS) reach streams not only
over the land surface but through groundwater discharge --- yet an in-stream sample cannot separate the
two, and watershed models have carried PFAS only from the surface into the subsurface, leaving the
groundwater contribution to in-stream PFAS, the quantity source attribution depends on, unquantified.
We advance watershed-scale PFAS fate and transport by resolving the two-way exchange between surface
water and groundwater: a process-based three-phase soil model (aqueous, Freundlich solid-phase, and
air--water interfacial sorption) carries PFAS from the land to the channel, and a coupled MODFLOW~6
groundwater flow-and-transport model returns the groundwater-borne PFAS to the stream in one continuous
mass balance. On the Rogue River, Michigan, downstream of the Wolverine World Wide legacy source, the
model is validated against an unusually rich record --- $5{,}383$ water-table heads (Nash--Sutcliffe
efficiency $0.91$), hundreds of groundwater PFOS measurements, streambed pore water at the discharge
interface, and a multi-analyte in-stream network. Two independent lines of evidence --- a model-free
multi-analyte fingerprint decomposition and a joint surface/groundwater calibration that avoids
double-counting --- show that legacy groundwater discharge is a first-order, spatially localized
contributor to the lower-mainstem in-stream PFOS load. A systematic global sensitivity and uncertainty
analysis identifies the controls on how PFAS partitions between the surface and groundwater pathways.
The result is a pathway-resolved, mass-conserving basis for attributing in-stream PFAS to its sources.
\end{abstract}
